% ------------------------ LaTeX Project Note Template ------------------------

% -- Set document class and font size ---
\documentclass[a4paper, 11pt]{article}

% -- Package imports --------------------
\usepackage[dvipsnames]{xcolor}
\usepackage{
    amsmath, amsthm, amssymb, mathtools, dsfont, units, physics,
    graphicx, wrapfig, float, booktabs, changepage,
    listings, inconsolata, pythonhighlight, subcaption,
    fancyhdr, hyperref, enumerate, enumitem, float }

% Load minted separately -- requires -shell-escape
\usepackage[most, minted]{tcolorbox}

% -- Page layout settings ---------------

% Set page margins
\usepackage[margin=3cm, a4paper, footskip=12mm]{geometry}

% Anchor footnotes to the bottom of the page
\usepackage[bottom]{footmisc}

% Set line spacing
\renewcommand{\baselinestretch}{1.2}

% Set spacing between paragraphs
\setlength{\parskip}{1.3mm}

% Allow multi-line equations to break onto the next page
\allowdisplaybreaks

% Get rid of paragraph indentation
\setlength{\parindent}{0pt}


% -- Page formatting settings -----------

% Set image captions to be italicized
\usepackage[font={it,footnotesize}]{caption}

% Set link colors for labeled items (blue), citations (red), URLs (orange)
\hypersetup{colorlinks=true, 
    linkcolor=RoyalBlue, 
    citecolor=RedOrange,
    urlcolor=ForestGreen
}

% Set font size for section titles (\large) and subtitles (\normalsize)
\usepackage{titlesec}
\titleformat{\section}{\Large\bfseries}{{\fontsize{19}{19}\selectfont\textreferencemark}\;\; }{0em}{}
\titleformat{\subsection}{\large\bfseries\selectfont}{\thesubsection\;\;\;}{0em}{}

% -- Table of Contents settings ---------
\usepackage[subfigure]{tocloft}

% Reduce spacing between sections in table of contents
\setlength{\cftbeforesecskip}{0.9ex}
\cftsetindents{section}{0em}{0em}
\renewcommand{\cfttoctitlefont}{\Large\bfseries}

% SAFE WAY: Remove section numbering width constraints in TOC 
% instead of opening unclosed \lrbox structures
\setcounter{tocdepth}{3} 
\renewcommand{\cftsecnumwidth}{0pt} 
\renewcommand{\cftsecindent}{0pt}

% -- Set path for images ----------------
\graphicspath{{figures/}{../figures/}}

% -- Math/Statistics commands -----------

% Add a reference number to a single line of a multi-line equation.
% Usage: "\numberthis\label{labelNameHere}" in an align or gather environment.
\newcommand\numberthis{\addtocounter{equation}{1}\tag{\theequation}}

% Set up Tikz for diagrams
\usepackage{tikz}
\usepackage{tikz-cd}
\usetikzlibrary{intersections, angles, quotes, calc, positioning}
\usetikzlibrary{arrows.meta}
\usepackage{pgfplots}
\pgfplotsset{compat=1.13}

\tikzset{
    force/.style={thick, {Circle[length=2pt]}-stealth, shorten <=-1pt}
}

% Set up different styles for theorems/remarks/notes/proofs/etc.
\usepackage{thmtools}
\usepackage[framemethod=TikZ]{mdframed}
\mdfsetup{skipabove=1em,skipbelow=0em}


\theoremstyle{definition}

\declaretheoremstyle[
    headfont=\bfseries\sffamily\color{ForestGreen!70!black}, bodyfont=\normalfont,
    mdframed={
        linewidth=2pt,
        rightline=false, topline=false, bottomline=false,
        linecolor=ForestGreen, backgroundcolor=ForestGreen!5,
    }
]{thmgreenbox}

\declaretheoremstyle[
    headfont=\bfseries\sffamily\color{NavyBlue!70!black}, bodyfont=\normalfont,
    mdframed={
        linewidth=2pt,
        rightline=false, topline=false, bottomline=false,
        linecolor=NavyBlue, backgroundcolor=NavyBlue!5,
    }
]{thmbluebox}

\declaretheoremstyle[
    headfont=\bfseries\sffamily\color{NavyBlue!70!black}, bodyfont=\normalfont,
    mdframed={
        linewidth=2pt,
        rightline=false, topline=false, bottomline=false,
        linecolor=NavyBlue
    }
]{thmblueline}

\declaretheoremstyle[
    headfont=\bfseries\sffamily\color{RawSienna!70!black}, bodyfont=\normalfont,
    mdframed={
        linewidth=2pt,
        rightline=false, topline=false, bottomline=false,
        linecolor=RawSienna, backgroundcolor=RawSienna!5,
    }
]{thmredbox}

\declaretheoremstyle[
    headfont=\bfseries\sffamily\color{RawSienna!70!black}, bodyfont=\normalfont,
    numbered=no,
    mdframed={
        linewidth=2pt,
        rightline=false, topline=false, bottomline=false,
        linecolor=RawSienna, backgroundcolor=RawSienna!1,
    },
    qed=\qedsymbol
]{thmproofbox}

\declaretheoremstyle[
    headfont=\bfseries\sffamily\color{NavyBlue!70!black}, bodyfont=\normalfont,
    numbered=no,
    mdframed={
        linewidth=2pt,
        rightline=false, topline=false, bottomline=false,
        linecolor=NavyBlue, backgroundcolor=NavyBlue!1,
    },
]{thmexplanationbox}

\declaretheorem[style=thmgreenbox, name=Key Idea]{definition}
\declaretheorem[style=thmbluebox, numbered=no, name=Example]{eg}
\declaretheorem[style=thmredbox, name=Claim]{prop}
\declaretheorem[style=thmredbox, name=Theorem]{theorem}

\declaretheorem[style=thmproofbox, name=Proof]{replacementproof}
% FIXED: Braces are now perfectly balanced
\renewenvironment{proof}[1][\proofname]{\vspace{-10pt}\begin{replacementproof}}{\end{replacementproof}}

\declaretheorem[style=thmexplanationbox, name=Explanation]{tmpexplanation}
\newenvironment{explanation}[1][]{\vspace{-10pt}\begin{tmpexplanation}}{\end{tmpexplanation}}


\declaretheorem[style=thmbluebox, numbered=no, name=Remark]{remark}
\declaretheorem[style=thmbluebox, numbered=no, name=Note]{note}

\newtheorem*{notation}{Notation}
\newtheorem*{previouslyseen}{As previously seen}
\newtheorem*{problem}{Problem}
\newtheorem*{observe}{Observe}
\newtheorem*{property}{Property}
\newtheorem*{intuition}{Intuition}

\newtcblisting{cppcode}{
    listing engine=minted,
    minted language=cpp,
    minted style=friendly,
    minted options={fontsize=\small, breaklines, autogobble},
    colback=gray!20,
    colframe=black,
    listing only,
    left=4mm,
    top=1mm,
    bottom=1mm,
}

\newtcblisting{bashcode}{
    listing engine=minted,
    minted language=bash,
    minted style=friendly,
    minted options={fontsize=\small, breaklines, autogobble},
    colback=gray!20,
    colframe=black,
    listing only,
    left=1mm,
    top=1mm,
    bottom=1mm,
}

\newtcblisting{tomlcode}{
    listing engine=minted,
    minted language=toml,
    minted style=friendly,
    minted options={fontsize=\small, breaklines, autogobble},
    colback=gray!20,
    colframe=black,
    listing only,
    left=4mm,
    top=1mm,
    bottom=1mm,
}

\newtcblisting{pythoncode}{
    listing engine=minted,
    minted language=python,
    minted style=friendly,
    minted options={fontsize=\small, breaklines, autogobble},
    colback=gray!20,
    colframe=black,
    listing only,
    left=4mm,
    top=1mm,
    bottom=1mm,
}


\newtcblisting{inputdeck}{
    listing engine=minted,
    minted language=text,
    minted style=friendly,
    minted options={fontsize=\small, breaklines, autogobble},
    colback=gray!20,
    colframe=black,
    listing only,
    left=4mm,
    top=1mm,
    bottom=1mm,
}

% -- Left/right header text -------------
\usepackage{fancyhdr}
\pagestyle{fancy}
\fancyhf{} % Explicitly reset everything safely

\setlength{\headheight}{15pt}
\renewcommand{\footrulewidth}{0pt}
\renewcommand{\headrulewidth}{0.4pt}

% Use the native automated marking engine safely
\renewcommand{\sectionmark}[1]{\markright{#1}}

\fancyhead[R]{\small\textit{\nouppercase{\rightmark}}}
\fancyfoot[C]{\thepage}